
\section{\points[5] Miscellaneous}
\subsection{\points[1] Memes}
\begin{enumerate}
\item \points[1] Draw a meme, write a joke, or share a fun observation related to course material.
\\~\\
This is a bonus question. You will receive full marks for any good-faith attempt.
\end{enumerate}

\clearpage

%% ============================================================
%% ROB 311 — Final Examination  (problems only, no preamble)
%% ============================================================

%% Points macro — matches quiz style \points[X]
%% Add to preamble if not already present:
%%   \newcommand{\points}[1]{\textbf{[#1 pts]}}


\subsection{\points[2] Evolutionary Learning}
For each question, circle the single best answer.
\begin{enumerate}

%% ---- MC1: Replicator dynamics ----
\item \points[1] In replicator dynamics, the fraction of the population playing strategy $i$
grows when:

\choicecorrect[The payoff to strategy $i$ exceeds the population-average payoff.] \\
\choice[Strategy $i$ is a strict Nash equilibrium of the underlying normal-form game.] \\
\choice[The population fraction of strategy $i$ exceeds $1/2$.] \\
\choice[All other strategies have strictly lower population fractions than strategy $i$.]

%% ---- MC2: Replicator dynamics / ESS ----
\item  \points[1] Which of the following correctly describes the relationship between
evolutionary stable strategies (ESS) and rest points of the replicator dynamics?

\choice[Every rest point of the replicator dynamics is an ESS.] \\
\choicecorrect[Every ESS is a rest point of the replicator dynamics, but not every rest point is an ESS.] \\
\choice[ESS and rest points of the replicator dynamics are exactly the same set.] \\
\choice[A mixed Nash equilibrium is always an ESS but never a rest point of the replicator dynamics.]
\end{enumerate}

\subsection{\points[2] Learning Theory}
For each question, circle the single best answer.

\begin{enumerate}

%% ---- MC3: PAC learning ----
\item  \points[1] In the PAC learning framework, a hypothesis $h$ is said to be
\textit{probably approximately correct} with parameters $(\varepsilon, \delta)$ if:

\choice[With probability at least $1 - \varepsilon$, the true error of $h$ is at most $\delta$.] \\
\choicecorrect[With probability at least $1 - \delta$, the true error of $h$ is at most $\varepsilon$.] \\
\choice[The empirical error of $h$ on the training set is at most $\varepsilon \cdot \delta$.] \\
\choice[The training set size exceeds $1/(\varepsilon \delta)$ examples.]

%% ---- MC4: Hoeffding's inequality ----
\item  \points[1] In Hoeffding's inequality, holding $\varepsilon$ fixed, doubling $n$:

\choice[Halves the bound on the deviation probability.] \\
\choice[Reduces the bound on the deviation probability by a factor of $e^{-2}$.] \\
\choicecorrect[Reduces the bound on the deviation probability by a factor of $e^{-2n\varepsilon^2}$, evaluated at the new $n$.] \\
\choice[Has no effect, since the bound depends only on $\varepsilon$ and the range of the random variables.]

\end{enumerate}

\clearpage

\section{\points[12] Search Problems}

\subsection{\points[12] Algorithm Tracing}

Consider the search problem over the graph below. The cost of each transition is represented by the weight of its associated edge. The heuristic value of each sate is represented in parentheses beside the associated vertex. The initial state is $s_0 = A$ and the goal set is $\mathcal{G} = \{G\}$.

\begin{figure}[H]
\centering
\begin{tikzpicture}[
  auto, semithick,
  every state/.style={draw=black, thick, fill=white, minimum size=7mm}
]
\node[state] (A) at (0, 0)    {$A$};
\node[state] (B) at (-3,-1.5) {$B$};
\node[state] (C) at ( 3,-1.5) {$C$};
\node[state] (D) at (-3,-3.5) {$D$};
\node[state] (E) at ( 3,-3.5) {$E$};
\node[state] (F) at ( 0,-3.5) {$F$};
\node[state] (G) at ( 0,-5.5) {$G$};

\path[->] (A) edge node[above left]  {$3$} (B);
\path[->] (A) edge node[above right] {$1$} (C);
\path[->] (B) edge node[left]  {$2$} (D);
\path[->] (C) edge[bend right] node[above] {$4$} (D);
\path[->] (C) edge node[right] {$1$} (E);
\path[->] (D) edge[bend right] node[below] {$1$} (F);
\path[->] (E) edge node[below] {$3$} (F);
\path[->, bend right=30] (F) edge node[above] {$2$} (D);
\path[->] (D) edge[bend right=15] node[below left] {$5$} (G);
\path[->] (F) edge node[right] {$2$} (G);

\node[above=1mm of A] {$(7)$};
\node[left=1mm  of B] {$(5)$};
\node[right=1mm of C] {$(8)$};
\node[left=1mm  of D] {$(3)$};
\node[right=1mm of E] {$(5)$};
\node[above=1mm of F] {$(2)$};
\node[right=1mm of G] {$(0)$};
\end{tikzpicture}
\end{figure}

\medskip\noindent
For each algorithm, provide the returned path, $p$ as a sequence of states, \texttt{<$\cdot$,$\cdot$,$\cdot$,$\dots$>} and its associated cost, $\textsf{cst}(p)$, or check the appropriate option.
\paragraph{Assumptions:}
\begin{itemize}
  \item when breaking ties, prefer the node whose path would come first alphabetically
  \item cycle checking (where applicable) is performed at \emph{export time} (before a
        successor is added to the open set).
\end{itemize}

\clearpage
\begin{enumerate}

\item \points[2] Breadth-First Search (no cycle checking). \\[5mm]
$p = \left\langle \ansbox[8cm]{A,B,D,G} \right\rangle, \textsf{cst}(p) = \ansbox[3cm]{10}$ \\[5mm]
\choice[The algorithm terminates but does not return a solution.]\\[5mm]
\choice[This algorithm will result in a cycle.]
\item \points[2] Breadth-First Search with inter-path cycle checking. \\[5mm]
$p = \left\langle \ansbox[8cm]{A,B,D,G} \right\rangle, \textsf{cst}(p) = \ansbox[3cm]{10}$ \\[5mm]
\choice[The algorithm terminates but does not return a solution.]\\[5mm]
\choice[This algorithm will result in a cycle.]
\item \points[2] Depth-First Search (no cycle checking). \\[5mm]
$p = \left\langle \ansbox[8cm]{A,B,D,F,D,F,...} \right\rangle, \textsf{cst}(p) = \numbox[3cm]$ \\[5mm]
\choice[The algorithm terminates but does not return a solution.]\\[5mm]
\choicecorrect[This algorithm will result in a cycle.]
\clearpage
\item \points[2] Depth-First Search with intra-path cycle checking. \\[5mm]
$p = \left\langle \ansbox[8cm]{A,B,D,F,G} \right\rangle, \textsf{cst}(p) = \ansbox[3cm]{8}$ \\[5mm]
\choice[The algorithm terminates but does not return a solution.]\\[5mm]
\choice[This algorithm will result in a cycle.]
\item \points[2] Cheapest-First (Uniform-Cost) Search. \\[5mm]
$p = \left\langle \ansbox[8cm]{A,C,E,F,G} \right\rangle, \textsf{cst}(p) = \ansbox[3cm]{7}$ \\[5mm]
\choice[The algorithm terminates but does not return a solution.]\\[5mm]
\choice[This algorithm will result in a cycle.]
\item \points[2] A*, using the heuristic values shown. \\[5mm]
$p = \left\langle \ansbox[8cm]{A,B,D,F,G} \right\rangle, \textsf{cst}(p) = \ansbox[3cm]{8}$ \\[5mm]
\choice[The algorithm terminates but does not return a solution.]\\[5mm]
\choice[This algorithm will result in a cycle.]

\end{enumerate}
\clearpage

%% ============================================================
\section{\points[16] Probabilistic Inference}

\subsection{\points[8] Dependence Separation}

Consider the Bayesian network below.

\begin{figure}[H]
\centering
\begin{tikzpicture}[
  auto, semithick,
  every state/.style={draw=black, thick, fill=white, minimum size=7mm},
  node distance=1.8cm
]
\node[state] (X) at (-1, 1.5) {$X$};
\node[state] (Y) at ( 1, 1.5) {$Y$};
\node[state] (Z) at ( 0, 0)   {$Z$};
\node[state] (W) at ( 0,-1.5) {$W$};
\node[state] (V) at ( 0,-3.0) {$V$};
\node[state] (U) at (-2, 0)   {$U$};

\path[-latex] (X) edge (Z);
\path[-latex] (Y) edge (Z);
\path[-latex] (Z) edge (W);
\path[-latex] (W) edge (V);
\path[-latex] (X) edge (U);
\end{tikzpicture}
\end{figure}

\noindent
Determine whether each statement is true or false, or indicate that there isn't enough information to determine the answer.

\begin{enumerate}

\item \points[2] $X \perp Y$ \\[5mm]
\choicecorrect[True]\choice[False]\choice[Not enough information]

\item \points[2] $X \perp Y \mid Z$ \\[5mm]
\choice[True]\choicecorrect[False]\choice[Not enough information]

\item \points[2] $X \perp Y \mid W$\\[5mm]
\choice[True]\choicecorrect[False]\choice[Not enough information]

\item \points[2] $U \perp Y \mid X$\\[5mm]
\choicecorrect[True]\choice[False]\choice[Not enough information]
\end{enumerate}
\begin{solution}
If students answered "Not enough information" instead of "False", only 1 mark penalty.
\end{solution}
\clearpage

\subsection{\points[8] Exact Inference via Bayes' Rule}

Consider the Bayesian network below, where $A, B, C, D$ are binary random variables over
$\{0,1\}$.

\begin{figure}[H]
\centering
\begin{tikzpicture}[
  auto, semithick,
  every state/.style={draw=black, thick, fill=white, minimum size=7mm}
]
\node[state] (A) at ( 0,  0)   {$A$};
\node[state] (B) at (-2, -1.5) {$B$};
\node[state] (C) at (-4, -3)   {$C$};
\node[state] (D) at ( 2, -3)   {$D$};

\path[-latex] (A) edge (B);
\path[-latex] (B) edge (C);
\path[-latex] (A) edge (D);
\path[-latex] (C) edge (D);
\end{tikzpicture}
\end{figure}

\noindent The conditional probability tables are:

\begin{table}[H]
\centering
\begin{tabular}{|c|c|}
\hline
$A$ & $P(A=1)$ \\\hline
    & $0.4$    \\\hline
\end{tabular}
\hspace{8mm}
\begin{tabular}{|c|c|}
\hline
$A$ & $P(B=1\mid A)$ \\\hline
$0$ & $0.3$          \\\hline
$1$ & $0.8$          \\\hline
\end{tabular}
\hspace{8mm}
\begin{tabular}{|c|c|}
\hline
$B$ & $P(C=1\mid B)$ \\\hline
$0$ & $0.2$          \\\hline
$1$ & $0.7$          \\\hline
\end{tabular}
\end{table}

\begin{table}[H]
\centering
\begin{tabular}{|c|c|c|}
\hline
$A$ & $C$ & $P(D=1\mid A,C)$ \\\hline
$0$ & $0$ & $0.10$           \\\hline
$0$ & $1$ & $0.50$           \\\hline
$1$ & $0$ & $0.60$           \\\hline
$1$ & $1$ & $0.90$           \\\hline
\end{tabular}
\end{table}

\begin{enumerate}
\item \points[8] Find the exact value of $P(A=1 \mid D=1)$. Provide your answer in fractional form.  Show your work. \\[5mm]
$\textsf{pr}(A=1|D=1) = \ansbox[3cm]{13} \; \Bigg/ \; \ansbox[3cm]{19}$\vspace{180pt}

\begin{solution}
By Bayes' rule:
\[
  P(A=1\mid D=1) = \frac{P(D=1\mid A=1)\,P(A=1)}{P(D=1)}.
\]
We marginalize over $B$ to compute $P(C=1\mid A)$:
\begin{align*}
  P(C=1\mid A=0) &= P(C=1\mid B=0)\,P(B=0\mid A=0) + P(C=1\mid B=1)\,P(B=1\mid A=0) \\
                 &= 0.2 \times 0.7 + 0.7 \times 0.3 = 0.14 + 0.21 = 0.35,\\[4pt]
  P(C=1\mid A=1) &= P(C=1\mid B=0)\,P(B=0\mid A=1) + P(C=1\mid B=1)\,P(B=1\mid A=1) \\
                 &= 0.2 \times 0.2 + 0.7 \times 0.8 = 0.04 + 0.56 = 0.60.
\end{align*}
Hence:
\[
  P(C=0\mid A=0)=1-0.35=0.65,\qquad P(C=0\mid A=1)=1-0.60=0.40.
\]
Now marginalize over $C$ to compute $P(D=1\mid A)$:
\begin{align*}
  P(D=1\mid A=0) &= P(D=1\mid A=0,C=0)\,P(C=0\mid A=0) \\
  &+ P(D=1\mid A=0,C=1)\,P(C=1\mid A=0) \\
                 &= 0.10 \times 0.65 + 0.50 \times 0.35 = 0.065 + 0.175 = 0.24,\\[4pt]
  P(D=1\mid A=1) &= P(D=1\mid A=1,C=0)\,P(C=0\mid A=1) \\
  &+ P(D=1\mid A=1,C=1)\,P(C=1\mid A=1) \\
                 &= 0.60 \times 0.40 + 0.90 \times 0.60 = 0.24 + 0.54 = 0.78.
\end{align*}
Then:
\[
  P(D=1) = P(D=1\mid A=0)\,P(A=0) + P(D=1\mid A=1)\,P(A=1)
         = 0.24\times 0.6 + 0.78\times 0.4
\]
\[
  = 0.144 + 0.312 = 0.456.
\]
Therefore:
\[
  P(A=1\mid D=1) = \frac{0.78 \times 0.4}{0.456} = \frac{0.312}{0.456} = \frac{13}{19} \approx 0.6842.
\]
\end{solution}
\end{enumerate}

\clearpage

\begin{center}
    \textit{(extra space)}
\end{center}

\clearpage

%% ============================================================
\section{\points[18] Markov Decision Processes}

\subsection{\points[12] Q-Value Iteration}

Consider an MDP with $\mathcal{S} = \{s_1, s_2, s_3\}$ and
$\mathcal{A} = \{\texttt{A},\texttt{B},\texttt{C}\}$.
Transitions and rewards are shown in the diagram below.

\begin{figure}[H]
\vspace{15pt}
\centering
\begin{tikzpicture}[
  auto, semithick,
  every state/.style={draw=black, thick, fill=white, minimum size=10mm}
]

%======================
% States
%======================
\node[state] (s1) at (0,3.5)  {$s_1$};
\node[state] (s2) at (6,3.5)  {$s_2$};
\node[state] (s3) at (12,3.5) {$s_3$};

%======================
% Actions
%======================
\node[state, dashed] (a1A) at (-1,0) {\texttt{A}};
\node[state, dashed] (a1B) at (0,0)  {\texttt{B}};
\node[state, dashed] (a1C) at (1,0)  {\texttt{C}};

\node[state, dashed] (a2A) at (5,7)  {\texttt{A}};
\node[state, dashed] (a2B) at (6,7)  {\texttt{B}};
\node[state, dashed] (a2C) at (7,7)  {\texttt{C}};

\node[state, dashed] (a3A) at (11,0) {\texttt{A}};
\node[state, dashed] (a3B) at (12,0) {\texttt{B}};
\node[state, dashed] (a3C) at (13,0) {\texttt{C}};

%======================
% State -> Action
%======================
\path[-latex, dashed] (s1) edge (a1A);
\path[-latex, dashed] (s1) edge (a1B);
\path[-latex, dashed] (s1) edge (a1C);

\path[-latex, dashed] (s2) edge (a2A);
\path[-latex, dashed] (s2) edge (a2B);
\path[-latex, dashed] (s2) edge (a2C);

\path[-latex, dashed] (s3) edge (a3A);
\path[-latex, dashed] (s3) edge (a3B);
\path[-latex, dashed] (s3) edge (a3C);

%======================
% Transitions
%======================

% s1
\path[-latex, overlay] (a1A) edge[out=110, in=100, looseness=1.4]
  node[pos=0.5, font=\small, fill=white, above=2mm] {$-1$} (s3);
\path[-latex] (a1B) edge[bend left=30] node[pos=0.5, font=\small, fill=white] {$+3$} (s2);
\path[-latex] (a1C) edge[bend left=10]  node[pos=0.5, font=\small, fill=white] {$+1$} (s2);
\path[-latex] (a1C) edge[bend right=30] node[above=5mm, pos=0.3, font=\small, fill=white] {$+1$} (s3);
% s2
\path[-latex] (a2A) edge[bend right=20] node[pos=0.5, font=\small, fill=white] {$+2$} (s1);
\path[-latex] (a2B) edge[bend right=10] node[pos=0.5, font=\small, fill=white] {$+5$} (s3);
\path[-latex] (a2C) edge[bend left=10]  node[pos=0.5, font=\small, fill=white] {$+2$} (s1);
\path[-latex] (a2C) edge[bend left=20]  node[pos=0.5, font=\small, fill=white] {$+2$} (s3);
% s3
\path[-latex] (a3A) edge[bend left=30]  node[below=5mm, pos=0.5, font=\small, fill=white] {$0$} (s1);
\path[-latex] (a3A) edge[bend right=10] node[pos=0.5, font=\small, fill=white] {$0$} (s2);
\path[-latex] (a3B) edge[loop below] node[pos=0.5, font=\small, fill=white] {$+4$} (a3B);
\path[-latex] (a3C) edge[bend right=30] node[pos=0.5, font=\small, fill=white] {$+1$} (s2);
\end{tikzpicture}
\end{figure}

\paragraph{Assumptions.}
\begin{itemize}
  \item $p(s'\mid s,a)$ is uniform over all successors shown from the action node.
  \item Discount factor $\gamma = 0.9$; time horizon $T = 2$.
  \item Initial Q-values: $q(s,a;\,2) \equiv 0$ for all $s \in \mathcal{S}$, $a \in \mathcal{A}$.
  \item All Q-values are updated synchronously via
  \[
    q(s,a;\,t) = \sum_{s'} p(s'\mid s,a)\Bigl[r(s,a,s') + \gamma\max_{a'}q(s',a';\,t+1)\Bigr].
  \]
\end{itemize}

\begin{enumerate}

\clearpage
\item \points[8] Fill in the Q-value tables below.


\renewcommand{\arraystretch}{1.9}
\begin{table}[H]
\centering
\begin{tabular}{|c|c|c|c|c|c|}
\hline
$s$ & $a$ & $q(s,a;\,2)$ & $q(s,a;\,1)$ & $q(s,a;\,0)$ & $\pi^*(a|s)$\\\hline
$s_1$ & \texttt{A} & $0$ & \ansnt{$-1$}   & \ansnt{$2.60$} & \ansnt{$0$} \\\hline
$s_1$ & \texttt{B} & $0$ & \ansnt{$3$}    & \ansnt{$7.50$} & \ansnt{$1$} \\\hline
$s_1$ & \texttt{C} & $0$ & \ansnt{$0.5$}  & \ansnt{$4.55$} & \ansnt{$0$} \\\hline
$s_2$ & \texttt{A} & $0$ & \ansnt{$2$}    & \ansnt{$4.70$} & \ansnt{$0$} \\\hline
$s_2$ & \texttt{B} & $0$ & \ansnt{$5$}    & \ansnt{$8.60$} & \ansnt{$1$} \\\hline
$s_2$ & \texttt{C} & $0$ & \ansnt{$2$}    & \ansnt{$5.15$} & \ansnt{$0$} \\\hline
$s_3$ & \texttt{A} & $0$ & \ansnt{$0.5$}  & \ansnt{$4.10$} & \ansnt{$0$} \\\hline
$s_3$ & \texttt{B} & $0$ & \ansnt{$4$}    & \ansnt{$7.60$} & \ansnt{$1$} \\\hline
$s_3$ & \texttt{C} & $0$ & \ansnt{$1$}    & \ansnt{$5.50$} & \ansnt{$0$} \\\hline
\end{tabular}
\end{table}

\item \points[2] Identify the greedy policy $\pi(a\mid s)$ derived from $q(s,a;\,0)$
for each $s \in \mathcal{S}$ and provide it in the table above.
\begin{solution}
\texttt{B} for all three states.
\end{solution}
\item \points[2] Without further computation, explain what happens to the greedy policy as $\gamma \to 0$.
\begin{solution}
As $\gamma \to 0$, future rewards are ignored, so the greedy policy approaches the policy that chooses the action with the largest expected \emph{immediate} reward at each state. For this MDP, this still gives \texttt{B} for all three states.
\end{solution}
\vspace{60pt}

\end{enumerate}

\clearpage

\subsection{\points[6] State Estimation in Partially Observable MDPs}

Consider the same MDP structure with $\mathcal{S}=\{s_1,s_2,s_3\}$ and
$\mathcal{A}=\{\texttt{A},\texttt{B},\texttt{C}\}$.
The transition probabilities under action \texttt{B} are:
\[
p(s_2\mid s_1,\texttt{B})=1,\quad
p(s_2\mid s_2,\texttt{B})=0.3,\quad
p(s_3\mid s_2,\texttt{B})=0.7,\quad
p(s_3\mid s_3,\texttt{B})=1.
\]
The observation model is:
\[
m(o\mid s) = \begin{cases} 0.6 & o = s \\ 0.2 & o \neq s \end{cases}
\]
for each of the other two states.

\begin{enumerate}

\item \points[6] The prior belief is $b(s_1)=0.2$,\; $b(s_2)=0.3$,\; $b(s_3)=0.5$.
The agent takes action $\texttt{B}$ and receives observation $o = s_1$.
Compute the posterior belief $b(s \mid a=\texttt{B},\, o=s_1)$.

\renewcommand{\arraystretch}{2.4}
\setlength{\tabcolsep}{18pt}   % default is 6pt
\begin{table}[H]
\centering
\begin{tabular}{|c|c|c|c|}
\hline
$s$ & $s_1$ & $s_2$ & $s_3$ \\\hline
$b(s)$ & $0.2$ & $0.3$ & $0.5$ \\\hline
$b(s \mid a=\texttt{B},\, o=s_1)$ & 
\ansnt{$0$} & \ansnt{$0.29$ or $0.2$} & \ansnt{$0.71$ or $0.8$} \\\hline
\end{tabular}
\end{table}
\begin{solution}
\textbf{Step 1: Predict} (propagate belief through action $\texttt{B}$).

\begin{align*}
  b'(s_1) &= b(s_1)\,p(s_1\mid s_1,\texttt{B}) + b(s_2)\,p(s_1\mid s_2,\texttt{B}) + b(s_3)\,p(s_1\mid s_3,\texttt{B}) \\
           &= 0 + 0 + 0 = 0, \\[4pt]
  b'(s_2) &= b(s_1)\,p(s_2\mid s_1,\texttt{B}) + b(s_2)\,p(s_2\mid s_2,\texttt{B}) + b(s_3)\,p(s_2\mid s_3,\texttt{B}) \\
           &= 0.2\times 1 + 0.3\times 0.3 + 0 = 0.2 + 0.09 = 0.29, \\[4pt]
  b'(s_3) &= b(s_1)\,p(s_3\mid s_1,\texttt{B}) + b(s_2)\,p(s_3\mid s_2,\texttt{B}) + b(s_3)\,p(s_3\mid s_3,\texttt{B}) \\
           &= 0 + 0.3\times 0.7 + 0.5\times 1 = 0.21 + 0.5 = 0.71.
\end{align*}

\medskip
\noindent\textbf{Step 2: Update} (condition on observation $o = s_1$).

\[
  b(s \mid a=\texttt{B}, o=s_1) \propto m(o=s_1 \mid s)\cdot b'(s).
\]

\begin{align*}
  m(s_1 \mid s_1)\cdot b'(s_1) &= 0.6 \times 0 = 0, \\
  m(s_1 \mid s_2)\cdot b'(s_2) &= 0.2 \times 0.29 = 0.058, \\
  m(s_1 \mid s_3)\cdot b'(s_3) &= 0.2 \times 0.71 = 0.142.
\end{align*}

Normaliser: $0 + 0.058 + 0.142 = 0.20$.

\[
  b(s_1 \mid a=\texttt{B}, o=s_1) = 0,\quad
  b(s_2 \mid a=\texttt{B}, o=s_1) = \frac{0.058}{0.2} = 0.29,\quad
  b(s_3 \mid a=\texttt{B}, o=s_1) = \frac{0.142}{0.2} = 0.71.
\]

Note that you can also skip the update step by seeing that the update shouldn't change the prediction probabilities. The observation is $o=s_1$. Since neither $s_2$ nor $s_3$ matches, they both get multiplied by the same likelihood $0.2$, so their relative proportions stay the same.

\end{solution}
\vspace{120pt}

\end{enumerate}

\clearpage

\begin{center}
    \textit{(extra space)}
\end{center}

\clearpage

%% ============================================================
\section{\points[20] Normal-Form Games}

\subsection{\points[8] Pure and Mixed Nash Equilibria}

Consider the following 2-player normal-form game.

\begin{table}[H]
\centering
\begin{tabular}{|c|c|c|}
\hline
              & $\texttt{L}$    & $\texttt{R}$    \\\hline
$\texttt{T}$  & $(3,\,1)$ & $(0,\,0)$ \\\hline
$\texttt{B}$  & $(0,\,0)$ & $(2,\,2)$ \\\hline
\end{tabular}
\end{table}

\begin{enumerate}

\item \points[3] Identify \emph{all} pure Nash equilibria, or state that none exist.
\\~\\
\checkboxcorrect[(\texttt{L},\texttt{T})]\checkbox[(\texttt{L},\texttt{B})]\checkbox[(\texttt{R},\texttt{T})]\checkboxcorrect[(\texttt{R},\texttt{B})]\choice[No pure equilibrium exists]
\vspace{100pt}

\item \points[5] Find the mixed Nash equilibrium strategy profile
\[
\pi = \bigl((p,\,1-p),\,(q,\,1-q)\bigr),
\]
where $p = \textsf{pr}[\texttt{T}]$ and $q = \textsf{pr}[\texttt{L}]$. Determine whether the mixed
equilibrium is \textbf{stable} or \textbf{unstable} under best-response dynamics, assuming
Agent~1 makes the first update.  Show your work.

\medskip

$p = \ansbox[3cm]{$\frac{2}{3}$}, \quad q = \ansbox[3cm]{$\frac{2}{5}$}$ \medskip \\
The equilibrium is \raisebox{1.75mm}{\choice[stable]} / \raisebox{1.75mm}{\choicecorrect[unstable]} under best-response play.
\begin{solution}
For the mixed Nash equilibrium, let $p=\textsf{pr}[\texttt{T}]$ and $q=\textsf{pr}[\texttt{L}]$.

Agent 1 is indifferent when
\[
3q = 2(1-q).
\]
Thus,
\[
3q = 2-2q \;\Rightarrow\; 5q=2 \;\Rightarrow\; q=\frac{2}{5}.
\]

Agent 2 is indifferent when
\[
1\cdot p = 2(1-p).
\]
Thus,
\[
p = 2-2p \;\Rightarrow\; 3p=2 \;\Rightarrow\; p=\frac{2}{3}.
\]

Therefore, the mixed Nash equilibrium is
\[
\pi = \left(\left(\frac{2}{3},\,\frac{1}{3}\right),\left(\frac{2}{5},\,\frac{3}{5}\right)\right).
\]

This mixed equilibrium is \textbf{unstable} under best-response dynamics. A small deviation causes best responses to push play toward one of the two pure equilibria rather than back to the mixed equilibrium.
\end{solution}
\vspace{120pt}

\end{enumerate}

\clearpage

\begin{center}
    \textit{(extra space)}
\end{center}

\clearpage

\subsection{\points[12] Subgame Perfect Equilibria}

Consider a 2-player, 2-stage normal-form game with the following payoff matrices.

\begin{table}[H]
\begin{minipage}{.5\linewidth}
\centering
\begin{tabular}{|c|c|c|}
\hline
             & $\texttt{c}$       & $\texttt{d}$       \\\hline
$\texttt{c}$ & $(2,\,2)$ & $(0,\,3)$ \\\hline
$\texttt{d}$ & $(3,\,0)$ & $(1,\,1)$ \\\hline
\end{tabular}
\caption*{Stage 1}
\end{minipage}%
\begin{minipage}{.5\linewidth}
\centering
\begin{tabular}{|c|c|c|}
\hline
             & $\texttt{e}$       & $\texttt{f}$       \\\hline
$\texttt{e}$ & $(5,\,5)$ & $(0,\,1)$ \\\hline
$\texttt{f}$ & $(1,\,0)$ & $(2,\,2)$ \\\hline
\end{tabular}
\caption*{Stage 2}
\end{minipage}
\end{table}

\noindent Consider the following strategy profile:
\begin{itemize}
  \item Stage 1: each agent plays $\texttt{c}$.
  \item Stage 2: each agent plays $\texttt{e}$ if the Stage-1 outcome was
        $(\texttt{c},\texttt{c})$; otherwise each agent plays $\texttt{f}$.
\end{itemize}

\begin{enumerate}

\item \points[4] Fill in the table below for each possible Stage-1 outcome ($\gamma = 1$).

\renewcommand{\arraystretch}{2.2}
\begin{table}[H]
\centering
\begin{tabular}{|c|c|c|c|c|}
\hline
Stage-1 outcome & Stage-1 rewards & Stage-2 outcome & Stage-2 rewards & Total returns \\\hline
$(\texttt{c},\texttt{c})$ & \ansnt{$(2,\,2)$} & \ansnt{$(\texttt{e},\texttt{e})$} & \ansnt{$(5,\,5)$} & \ansnt{$(7,\,7)$} \\\hline
$(\texttt{d},\texttt{d})$ & \ansnt{$(1,\,1)$} & \ansnt{$(\texttt{f},\texttt{f})$} & \ansnt{$(2,\,2)$} & \ansnt{$(3,\,3)$} \\\hline
$(\texttt{c},\texttt{d})$ & \ansnt{$(0,\,3)$} & \ansnt{$(\texttt{f},\texttt{f})$} & \ansnt{$(2,\,2)$} & \ansnt{$(2,\,5)$} \\\hline
$(\texttt{d},\texttt{c})$ & \ansnt{$(3,\,0)$} & \ansnt{$(\texttt{f},\texttt{f})$} & \ansnt{$(2,\,2)$} & \ansnt{$(5,\,2)$} \\\hline
\end{tabular}
\end{table}

\clearpage
\item \points[4] Is this strategy profile a Nash equilibrium of the 2-stage game?
Justify your answer.
\\~\\
\choicecorrect[Yes]\choice[No]
\begin{solution}

Given the specified continuation strategies, the total returns from Stage 1 are:
\[
(\texttt{c},\texttt{c}) \mapsto (7,7),\quad
(\texttt{d},\texttt{d}) \mapsto (3,3),\quad
(\texttt{c},\texttt{d}) \mapsto (2,5),\quad
(\texttt{d},\texttt{c}) \mapsto (5,2).
\]
Thus, if one player chooses $\texttt{c}$, the other prefers $\texttt{c}$ since $7>5$, and if one player chooses $\texttt{d}$, the other prefers $\texttt{d}$ since $3>2$. Therefore $(\texttt{c},\texttt{c})$ is a Nash equilibrium of the 2-stage game.
\end{solution}
\vspace{160pt}

\item \points[4] Is this strategy profile a subgame perfect equilibrium?
Justify your answer.
\\~\\
\choicecorrect[Yes]\choice[No]
\begin{solution}
After every possible Stage-1 outcome, the prescribed Stage-2 action profile is a Nash equilibrium of the Stage-2 game:
\[
(\texttt{e},\texttt{e}) \text{ is a Nash equilibrium, and } (\texttt{f},\texttt{f}) \text{ is also a Nash equilibrium.}
\]
Moreover, at Stage 1, given these continuation payoffs, $(\texttt{c},\texttt{c})$ is a Nash equilibrium. Hence the strategy profile is a subgame perfect equilibrium.
\end{solution}
\vspace{80pt}

\end{enumerate}

\clearpage

%% ============================================================
\section{\points[18] Extensive-Form Games}

\subsection{\points[8] $\alpha$/$\beta$-Pruning}

Consider the 4-ply zero-sum game tree below. The maximiser acts at the root and at every even depth; the minimiser acts at every odd depth. Leaf values are the maximiser's payoffs.
\begin{figure}[H]
\centering
\begin{tikzpicture}[auto, semithick,
    every state/.style={draw=black, thick, fill=white, minimum size=9mm},
    scale=0.92, transform shape]

\node[state] (A) at (0,0)      {$A$:\ansnt{4}};
\node[state] (B) at (-4.8,-1.8)  {$B$:\ansnt{3}};
\node[state] (C) at ( 4.8,-1.8)  {$C$:\ansnt{4}};
\node[state] (D) at (-7.2,-3.6)  {$D$:\ansnt{3}};
\node[state] (G) at (-2.4,-3.6)  {$G$:\ansnt{6}};
\node[state] (J) at ( 2.4,-3.6)  {$J$:\ansnt{4}};
\node[state] (M) at ( 7.2,-3.6)  {$M$:\ansnt{7}};
\node[state] (E) at (-8.4,-5.4){$E$:\ansnt{3}};
\node[state] (F) at ( -6.0,-5.4){$F$:\ansnt{1}};
\node[state] (H) at ( -3.6,-5.4){$H$:\ansnt{2}};
\node[state] (I) at ( -1.2,-5.4){$I$:\ansnt{6}};
\node[state] (K) at (  1.2,-5.4){$K$:\ansnt{1}};
\node[state] (L) at (  3.6,-5.4){$L$:\ansnt{4}};
\node[state] (N) at (  6.0,-5.4){$N$:\ansnt{7}};
\node[state] (O) at (  8.4,-5.4){$O$:\ansnt{X}};

\node[state, rectangle] (l1)  at (-9.0,-7.2) {\small $l_1{:}3$};
\node[state, rectangle] (l2)  at ( -7.8,-7.2) {\small $l_2{:}5$};
\node[state, rectangle] (l3)  at ( -6.6,-7.2) {\small $l_3{:}7$};
\node[state, rectangle] (l4)  at ( -5.4,-7.2) {\small $l_4{:}1$};
\node[state, rectangle] (l5)  at ( -4.2,-7.2) {\small $l_5{:}2$};
\node[state, rectangle] (l6)  at ( -3.0,-7.2) {\small $l_6{:}4$};
\node[state, rectangle] (l7)  at ( -1.8,-7.2) {\small $l_7{:}6$};
\node[state, rectangle] (l8)  at ( -0.6,-7.2) {\small $l_8{:}6$};
\node[state, rectangle] (l9)  at (  0.6,-7.2) {\small $l_9{:}4$};
\node[state, rectangle] (l10) at (  1.8,-7.2) {\small $l_{10}{:}1$};
\node[state, rectangle] (l11) at (  3.0,-7.2) {\small $l_{11}{:}5$};
\node[state, rectangle] (l12) at (  4.2,-7.2) {\small $l_{12}{:}4$};
\node[state, rectangle] (l13) at (  5.4,-7.2) {\small $l_{13}{:}8$};
\node[state, rectangle] (l14) at (  6.6,-7.2) {\small $l_{14}{:}7$};
\node[state, rectangle] (l15) at (  7.8,-7.2) {\small $l_{15}{:}9\ansnt{X}$};
\node[state, rectangle] (l16) at (  9.0,-7.2) {\small $l_{16}{:}8\ansnt{X}$};

\path (A) edge (B); \path (A) edge (C);
\path (B) edge (D); \path (B) edge (G);
\path (C) edge (J); \path (C) edge (M);
\path (D) edge (E); \path (D) edge (F);
\path (G) edge (H); \path (G) edge (I);
\path (J) edge (K); \path (J) edge (L);
\path (M) edge (N); \path (M) edge (O);
\path (E) edge (l1);  \path (E) edge (l2);
\path (F) edge (l3);  \path (F) edge (l4);
\path (H) edge (l5);  \path (H) edge (l6);
\path (I) edge (l7);  \path (I) edge (l8);
\path (K) edge (l9);  \path (K) edge (l10);
\path (L) edge (l11); \path (L) edge (l12);
\path (N) edge (l13); \path (N) edge (l14);
\path (O) edge (l15); \path (O) edge (l16);
\end{tikzpicture}
\end{figure}

\noindent
\begin{enumerate}
    \item \points[8] Apply $\alpha/\beta$-pruning (depth-first, left-to-right).
Fill in the backed-up value for every \emph{explored} node and
\textbf{cross out} every \emph{unexplored} node or leaf.
You do not need to record $\alpha$/$\beta$ values explicitly.
\begin{solution}
Depth structure: $A$ (max) $\to$ $B,C$ (min) $\to$ $D,G,J,M$ (max) $\to$ $E,F,H,I,K,L,N,O$ (min) $\to$ leaves.

\medskip
\noindent \textbf{Left subtree ($A \to B$):}

\smallskip
$D$ (max): $E=\min(3,5)=3$ and $F=\min(7,1)=1$, so $D=\max(3,1)=3$.

$B$ (min): current best is $3$. Now explore $G$.

$G$ (max): $H=\min(2,4)=2$ and $I=\min(6,6)=6$, so $G=\max(2,6)=6$.

Hence $B=\min(3,6)=3$. Therefore after the left subtree, $A \ge 3$.

\medskip
\noindent \textbf{Right subtree ($A \to C$):}

\smallskip
$J$ (max): $K=\min(4,1)=1$ and $L=\min(5,4)=4$, so $J=\max(1,4)=4$.

Thus $C$ (min) currently has value at most $4$.

$M$ (max): explore $N=\min(8,7)=7$. Then $M \ge 7$, which is already enough to make $C$ keep its earlier choice $J=4$. Therefore prune $O$ and its children $l_{15}, l_{16}$.

So $M=7$ for the search, and
\[
C=\min(4,7)=4.
\]

Finally,
\[
A=\max(3,4)=4.
\]

\medskip
\noindent \textbf{Summary of explored/pruned nodes:}
\begin{itemize}
  \item \textbf{Explored:} $A=4$, $B=3$, $C=4$, $D=3$, $G=6$, $J=4$, $M=7$, $E=3$, $F=1$, $H=2$, $I=6$, $K=1$, $L=4$, $N=7$.
  \item \textbf{Pruned:} $O$, $l_{15}$, and $l_{16}$.
\end{itemize}

\medskip
\noindent \textit{Note}: $O$, $l_{15}$, and $l_{16}$ are crossed out (unexplored). All other explored nodes and leaves are filled with their search-time backed-up values.
\end{solution}
\end{enumerate}



\clearpage

\subsection{\points[10] Monte-Carlo Tree Search}

The following sub-questions are \textbf{independent} of one another.

\begin{enumerate}

\item \points[5]
Consider a 2-player game in which both agents are learning a policy using
Monte-Carlo Tree Search (MCTS). The state of the game tree after some number of
iterations is shown below. Edge labels are $w/n$, where $n$ is the number of times
the action was selected and $w$ is the number of wins for the \emph{agent who took
that action}. Agents alternate turns by depth; Agent~1 starts at the root.

\begin{figure}[H]
\centering
\begin{tikzpicture}[auto, semithick,
    every state/.style={draw=black, thick, fill=white, minimum size=6mm}]
\node[state] at ( 0, 0)  (A) {$A$};
\node[state] at (-3,-2)  (B) {$B$};
\node[state] at ( 3,-2)  (C) {$C$};
\node[state] at (-4,-4)  (D) {$D$};
\node[state] at (-2,-4)  (E) {$E$};
\node[state] at (-3,-6)  (F) {$F$};
\node[state] at (-1,-6)  (G) {$G$};
\node[state] at (-2,-8)  (H) {$H$};
\node[state] at (0,-8)  (I) {$I$};
\node[state] at ( 2,-4)  (J) {$J$};
\node[state] at ( 4,-4)  (K) {$K$};

\path (A) edge node[midway, above left]  {$3/6$} (B);
\path (A) edge node[midway, above right] {$1/3$} (C);
\path (B) edge node[midway, above left]  {$2/3$} (D);
\path (B) edge node[midway, above right] {$1/3$} (E);
\path (E) edge node[midway, above left]  {$0/1$} (F);
\path (E) edge node[midway, above right] {$??$}  (G);
\path (G) edge node[midway, above left]  {$0/1$} (H);
\path (G) edge node[midway, above right] {$0/1$} (I);
\path (C) edge node[midway, above left]  {$1/2$} (J);
\path (C) edge node[midway, above right] {$0/1$} (K);
\end{tikzpicture}
\end{figure}

\noindent
During the next iteration, the selection phase follows the path
$\langle A, B, E, G \rangle$ and node $J'$ is expanded as a new child of $G$.
The simulation runs 6 transitions from $J'$, and the
\emph{first turn-taker} (i.e., the turn-taker at $A$) wins the game.

\medskip\noindent
Compute the updated $w/n$ values for the four edges on the selection path.

\begin{table}[H]
\centering
\renewcommand{\arraystretch}{2.2}
\begin{tabular}{|c|c|c|c|}
\hline
$q(A,\langle A,B\rangle)$ & $q(B,\langle B,E\rangle)$ & $q(E,\langle E,G\rangle)$ & $q(G,\langle G,J'\rangle)$ \\\hline
\ansnt{$4/7$} & \ansnt{$1/4$} & \ansnt{$3/3$} & \ansnt{$0/1$} \\\hline
\end{tabular}
\end{table}
\ansnt{Marking scheme: 1 for each box, +1 for recognizing that $q(E,\langle E,G\rangle)$ is out of 3. If student got Agent 2 wins but everything else is consistent (3/7, 2/4, 2/3, 1/1), only -2 marks.}
% \vspace{20pt}

%% ---- MCTS Q2: softmax selection ----
\item \points[5]
The game tree after 20 iterations is shown below. Edge labels are $w/n$,
where $w$ is the number of wins for the agent who took that action.
Agents alternate turns by depth; Agent~1 starts at the root.

\begin{figure}[H]
\centering
\begin{tikzpicture}[auto, semithick,
    every state/.style={draw=black, thick, fill=white, minimum size=6mm}]
\node[state] at ( 0, 0)   (A) {$A$};
\node[state] at (-4,-2)   (B) {$B$};
\node[state] at ( 0,-2)   (C) {$C$};
\node[state] at ( 4,-2)   (D) {$D$};
\node[state] at (-5,-4)   (E) {$E$};
\node[state] at (-3,-4)   (F) {$F$};
\node[state] at (-1,-4)   (G) {$G$};
\node[state] at ( 1,-4)   (H) {$H$};
\node[state] at ( 3,-4)   (I) {$I$};
\node[state] at ( 5,-4)   (J) {$J$};

\path (A) edge node[midway, above left]  {$4/8$}  (B);
\path (A) edge node[midway, left]        {$3/8$}  (C);
\path (A) edge node[midway, above right] {$2/4$}  (D);
\path (B) edge node[midway, above left]  {$3/5$}  (E);
\path (B) edge node[midway, above right] {$2/3$}  (F);
\path (C) edge node[midway, above left]  {$1/4$}  (G);
\path (C) edge node[midway, above right] {$2/4$}  (H);
\path (D) edge node[midway, above left]  {$1/2$}  (I);
\path (D) edge node[midway, above right] {$2/2$}  (J);
\end{tikzpicture}
\end{figure}

\noindent
Both players use softmax selection with temperature $\tau = 1$. At each
node, the probability of selecting child $i$ is
\[
  \pi(i) = \frac{\text{exp}\bigl(w_i/n_i\bigr)}{\displaystyle\sum_j \text{exp}\bigl(w_j/n_j\bigr)},
\]
where the sum is over all children of the current node.

\medskip\noindent
Compute the probability that node $G$ lies on the selection path during the
21\textsuperscript{st} iteration. Show all steps.

\medskip\noindent
The probability that $G$ is on the selection path during the 21\textsuperscript{st} iteration is \ansbox[2cm]{0.134}.
\begin{solution}
\[
P(G \text{ is on the path}) = P(C\mid A)\,P(G\mid C).
\]

At node $A$, the edge win-rates are
\[
\frac{4}{8}=0.5,\qquad \frac{3}{8}=0.375,\qquad \frac{2}{4}=0.5.
\]
Using softmax with $\tau=1$,
\[
P(C\mid A)=\frac{e^{3/8}}{e^{4/8}+e^{3/8}+e^{2/4}}
=\frac{e^{0.375}}{e^{0.5}+e^{0.375}+e^{0.5}}
\approx 0.3061.
\]

At node $C$, the edge win-rates are
\[
\frac{1}{4}=0.25,\qquad \frac{2}{4}=0.5.
\]
So
\[
P(G\mid C)=\frac{e^{1/4}}{e^{1/4}+e^{2/4}}
=\frac{e^{0.25}}{e^{0.25}+e^{0.5}}
\approx 0.4377.
\]

Hence,
\[
P(G \text{ is on the path}) \approx 0.3061 \times 0.4377 \approx 0.134.
\]
\end{solution}
\vspace{140pt}

\end{enumerate}

\clearpage

\begin{center}
    \textit{(extra space)}
\end{center}

\clearpage

%% ============================================================
\section{\points[14] Stochastic Games}

\subsection{\points[14] Nash Value Iteration}

Consider a 2-player stochastic game with states $\mathcal{S}=\{\texttt{s1},\texttt{s2}\}$
and joint actions,
\[\mathcal{A}=\{(\texttt{c,c}),(\texttt{c,d}),(\texttt{d,c}),(\texttt{d,d})\}\]

\medskip\noindent
The rewards that Agent 1 would receive for each state and joint action is shown below:

\begin{table}[H]
\centering
\begin{tabular}{|c|c|c|c|c|}
\hline
$s$ & \texttt{c,c} & \texttt{c,d} & \texttt{d,c} & \texttt{d,d} \\\hline
\texttt{s1} & $3$ & $0$ & $4$ & $1$ \\\hline
\texttt{s2} & $5$ & $1$ & $1$ & $5$ \\\hline
\end{tabular}
\end{table}

\noindent
\noindent The probability of transitioning to state \texttt{s2} given each state and joint action is shown below:

\begin{table}[H]
\centering
\begin{tabular}{|c|c|c|c|c|}
\hline
$s$ & \texttt{c,c} & \texttt{c,d} & \texttt{d,c} & \texttt{d,d} \\\hline
\texttt{s1} & $0.20$ & $0.50$ & $0.50$ & $0.90$ \\\hline
\texttt{s2} & $0.70$ & $0.40$ & $0.40$ & $0.80$ \\\hline
\end{tabular}
\end{table}

\noindent
\noindent Assume that $T=2$ and the agents use a discount factor $\gamma=1.0$.

\begin{enumerate}

\item \points[4]
Find the Nash equilibrium of the one-shot game at each state ($t=1$) and compute
$v_1^*(s;\,t=1)$.  Show your reasoning.
\begin{solution}
At $t=1$, we solve the one-shot game at each state.

For $\texttt{s1}$, the game is
\[
\begin{array}{c|cc}
 & \texttt{c} & \texttt{d}\\\hline
\texttt{c} & (3,3) & (0,4)\\
\texttt{d} & (4,0) & (1,1)
\end{array}
\]
For each player, $\texttt{d}$ strictly dominates $\texttt{c}$, so the unique Nash equilibrium is
\[
(\texttt{d},\texttt{d}),
\qquad
v_1^*(\texttt{s1};1)=1.
\]

For $\texttt{s2}$, the game is
\[
\begin{array}{c|cc}
 & \texttt{c} & \texttt{d}\\\hline
\texttt{c} & (5,5) & (1,1)\\
\texttt{d} & (1,1) & (5,5)
\end{array}
\]
The pure Nash equilibria are
\[
(\texttt{c},\texttt{c}) \quad \text{and} \quad (\texttt{d},\texttt{d}),
\]
and both give Agent 1 payoff $5$. Hence
\[
v_1^*(\texttt{s2};1)=5.
\]
\end{solution}
\vspace{120pt}

\clearpage
\item \points[8]
Compute $q_1^*(s,a;\,t=0)$ for all $(s,a)$ using
\[
  q_1^*(s,a;\,t=0) = r_1(s,a)
    + \gamma\sum_{s'} p(s'\mid s,a)\,v_1^*(s';\,t=1).
\]
Fill in the table. Show your work.

\renewcommand{\arraystretch}{2.3}
\begin{table}[H]
\centering
\begin{tabular}{|c|c|c|c|c|}
\hline
$s$ & \hspace{1cm}(\texttt{c,c})\hspace{1cm} & \hspace{1cm}(\texttt{c,d})\hspace{1cm} & \hspace{1cm}(\texttt{d,c})\hspace{1cm} & \hspace{1cm}(\texttt{d,d})\hspace{1cm} \\\hline
\texttt{s1} & \ansnt{$4.8$} & \ansnt{$3.0$} & \ansnt{$7.0$} & \ansnt{$5.6$} \\\hline
\texttt{s2} & \ansnt{$8.8$} & \ansnt{$3.6$} & \ansnt{$3.6$} & \ansnt{$9.2$} \\\hline
\end{tabular}
\end{table}
\begin{solution}
Using
\[
q_1^*(s,a;0)=r_1(s,a)+\sum_{s'} p(s'\mid s,a)\,v_1^*(s';1),
\]
with
\[
v_1^*(\texttt{s1};1)=1,\qquad v_1^*(\texttt{s2};1)=5,
\]
we get:

\[
q_1^*(\texttt{s1},\texttt{c,c};0)=3+0.8(1)+0.2(5)=4.8,
\]
\[
q_1^*(\texttt{s1},\texttt{c,d};0)=0+0.5(1)+0.5(5)=3.0,
\]
\[
q_1^*(\texttt{s1},\texttt{d,c};0)=4+0.5(1)+0.5(5)=7.0,
\]
\[
q_1^*(\texttt{s1},\texttt{d,d};0)=1+0.1(1)+0.9(5)=5.6,
\]

\[
q_1^*(\texttt{s2},\texttt{c,c};0)=5+0.3(1)+0.7(5)=8.8,
\]
\[
q_1^*(\texttt{s2},\texttt{c,d};0)=1+0.6(1)+0.4(5)=3.6,
\]
\[
q_1^*(\texttt{s2},\texttt{d,c};0)=1+0.6(1)+0.4(5)=3.6,
\]
\[
q_1^*(\texttt{s2},\texttt{d,d};0)=5+0.2(1)+0.8(5)=9.2.
\]

\end{solution}
\clearpage
\item \points[2]
Without computing a Nash equilibrium at $t=0$, identify which joint action at state
\texttt{s1} is \emph{never} part of a Nash equilibrium at $t=0$, and explain why in one
sentence.
\begin{solution}
At state \texttt{s1}, the only Nash equilibrium at $t=0$ is $(\texttt{d},\texttt{c})$. Therefore the joint actions
\[
\ansnt{(\texttt{c},\texttt{c}),\;(\texttt{c},\texttt{d}),\;\text{and}\;(\texttt{d},\texttt{d})}
\]
are never part of a Nash equilibrium at $t=0$.

Specifically:
\begin{itemize}
  \item At $(\texttt{c},\texttt{c})$, Agent 1 gets $4.8$ and can improve by switching to \texttt{d}, which gives $7.0$.
  \item At $(\texttt{c},\texttt{d})$, Agent 1 gets $3.0$ and can improve by switching to \texttt{d}, which gives $5.6$.
  \item At $(\texttt{d},\texttt{d})$, Agent 2 gets $5.6$ and can improve by switching to \texttt{c}, which gives $7.0$.
\end{itemize}
\end{solution}
\vspace{60pt}

\end{enumerate}

\clearpage

\begin{center}
    \textit{(extra space)}
\end{center}

\clearpage

%% ============================================================
\section{\points[10] Search Problems (Advanced)}
\subsection{\points[10] Modelling}

You are planning a train journey and want to minimize total travel time.
You start at station $A$ and want to reach station $F$.
However, there are a few constraints:
\begin{itemize}
  \item Some routes require express train tickets, and you want to avoid taking more
        than 2 express routes total.
  \item You must visit all scenic viewpoints before reaching $F$
        (at stations $B$, $C$, and $E$).
  \item Some train routes only run in one direction.
  \item Travel times between stations vary.
\end{itemize}

The available train routes are shown in the graph below (shaded nodes are scenic viewpoints):

\begin{figure}[H]
\centering
\begin{tikzpicture}[
  auto, node distance=3cm, semithick,
  every state/.style={draw=black, thick, fill=white, minimum size=10mm},
  viewpoint/.style={draw=black, thick, fill=gray!25, minimum size=10mm}
]
\node[state]     (a) at (0, 0)    {$A$};
\node[viewpoint] (b) at (-3,-2.5) {$B$};
\node[viewpoint] (c) at ( 3,-2.5) {$C$};
\node[state]     (d) at (-5,-5)   {$D$};
\node[viewpoint] (e) at ( 5,-5)   {$E$};
\node[state]     (f) at ( 0,-7)   {$F$};

\path[->]   (a) edge node[left,  xshift=-4mm, text width=1.5cm, align=center]
                       {30 min\\\textit{express}} (b);
\path[<->]  (a) edge node[right, xshift=4mm,  text width=1.5cm, align=center]
                       {45 min\\\textit{local}}   (c);
\path[<->]  (b) edge node[left,              text width=1.5cm, align=center]
                       {20 min\\\textit{local}}   (d);
\path[->]   (b) edge node[anchor=center, xshift=1cm, yshift=-1.25cm,
                           text width=1.5cm, align=center]
                       {35 min\\\textit{express}} (e);
\path[->]   (c) edge node[anchor=center, xshift=-1cm, yshift=-1.25cm,
                           text width=1.5cm, align=center]
                       {25 min\\\textit{express}} (d);
\path[<->]  (c) edge node[midway, text width=1.5cm, align=center]
                       {40 min\\\textit{local}}   (e);
\path[->]   (d) edge node[below left, text width=1.5cm, align=center]
                       {50 min\\\textit{local}}   (f);
\path[->]   (e) edge node[below right, text width=1.5cm, align=center]
                       {30 min\\\textit{local}}   (f);
\end{tikzpicture}
\end{figure}

\begin{enumerate}

\clearpage 
\item \points[2] Explain why a simple search (i.e., one that finds any path from $A$ to
$F$) over the graph above would be \emph{insufficient} for this problem.
Give a specific example.
\begin{solution}
A simple search over only stations is insufficient because the same station can represent different constraint satisfaction states. For example, reaching $E$ without having visited $B$ and $C$ violates the requirement, while reaching $E$ after visiting them is valid. Similarly, reaching a station after using 2 express routes prevents taking further express routes, while reaching it with fewer does not.
\end{solution}
\vspace{2cm}
\item \points[8] Model this as a search problem by filling in each component below. You may not need all of the boxes.
  \begin{enumerate}

  \item \points[2] Each state is a tuple,
  \[\left(\ansbox[2cm]{station},\ansbox[2cm]{$b$},\ansbox[2cm]{$c$},\ansbox[2cm]{$e$},\ansbox[2cm]{$m$} \right),\]
  where: \\[2mm]
  \begin{itemize}
      \item \; \ansnt{station: the current station (one of $A,B,C,D,E,F$)} \\[5mm]
      \item \; \ansnt{$b \in \{\text{T},\text{F}\}$: whether $B$ has been visited} \\~\\
      \item \; \ansnt{$c \in \{\text{T},\text{F}\}$: whether $C$ has been visited} \\~\\
      \item \; \ansnt{$e \in \{\text{T},\text{F}\}$: whether $E$ has been visited} \\~\\
      \item \; \ansnt{$m \in \{0,1,2\}$: number of express routes used so far}
  \end{itemize}
  \begin{solution}
  Marking scheme:
  \begin{itemize}
      \item 2 marks for minimal correct state: location, visited viewpoints (set or bit mask), express count
      \item 1.5 marks for slightly redundant but correct state. E.g., included time and path
      \item 0-1 marks for incomplete state or major conceptual confusion (e.g., relied on search algorithm to check for express routes taken)
  \end{itemize}
  \end{solution}
  \item \points[1] Write the initial state using your representation.
    \[s_0 = \left(\ansbox[2cm]{$A$},\ansbox[2cm]{F},\ansbox[2cm]{F},\ansbox[2cm]{F},\ansbox[2cm]{0}\right)\]
  \item \points[1] Describe which states are goal states using your representation.
  % \vspace{100pt}
  \begin{solution}
      $s_g$ = ($F$, T, T, T, any $m$)
  \end{solution}

  \item \points[2] For a generic state $s$, describe how to determine its successors
        (i.e., the states reachable via a single action).
  % \vspace{100pt}
  \begin{solution}
      Explore edges normally, setting the bit mask to T if it visits the corresponding viewpoint. Explore the express edges only if $m < 2$ and decrease express count by 1.
      \\~\\
      Marking scheme: only 1 mark if student only describes the station portion of successor, and not the complete state.
  \end{solution}

  \item \points[2] Describe an admissible heuristic for this problem.
        Explain in one or two sentences why your heuristic is admissible.
  % \vspace{80pt}
  \begin{solution}
  One helpful heuristic: Take the minimum travel time from the current station to $F$ in the graph while \textit{ignoring the viewpoint and express-ticket constraints}. This is admissible because removing constraints can only make the remaining trip cheaper, so it is a lower bound on the true remaining cost.
  \\~\\
  Any technically admissible heuristic is also accepted. E.g. $h(x)=0$ is not very helpful but technically admissible.
  \end{solution}

  \end{enumerate}

\end{enumerate}

\clearpage

\begin{center}
    \textit{(extra space)}
\end{center}

\clearpage

%% ============================================================
\section{\points[8] Markov Decision Processes (Advanced)}
\subsection{\points[8] Modelling}

A small autonomous rover operates in a $1 \times 3$ grid of cells:
\texttt{left}, \texttt{centre}, and \texttt{right}.
The rover has a battery level that is either \texttt{high} or \texttt{low}.

\medskip\noindent
At each time-step the rover chooses one of three actions:
\begin{itemize}
  \item \texttt{move}: attempt to move one cell to the right.
    \begin{itemize}
      \item if the battery is \texttt{high} this succeeds with probability $0.9$ (moves right) and
            fails with probability $0.1$ (stays in current cell)
      \item if battery is \texttt{low} this succeeds with probability $0.5$ and
            fails with probability $0.5$
      \item if the rover is already at \texttt{right}, the action always fails (stays)
    \end{itemize}
  \item \texttt{stay}: remain in the current cell (succeeds all the time and the battery will not change)
  \item \texttt{charge}: remain in the current cell and set battery to \texttt{high}
        (suceeds all the time but only available at \texttt{left})
\end{itemize}
\medskip\noindent
The rewards are as follows:
\begin{itemize}
  \item $+10$ for reaching \texttt{right} for the first time
  \item $-2$ for \texttt{charge} action (charging takes time)
  \item $-1$ for every other transition (as a small time penalty)
\end{itemize}

\medskip\noindent
The rover wants to reach \texttt{right} while maximising total (undiscounted) reward.

\begin{enumerate}

\item \points[1] Assume each state is represented as a tuple $(l,b)$, where $l \in \left\lbrace \texttt{left}, \texttt{center}, \texttt{right} \right\rbrace$ represents the location and $b \in \left\lbrace \texttt{low}, \texttt{high} \right\rbrace$ represents the battery lavel. Define the state space $\mathcal{S}$ for this MDP by providing the list of all the states.
\\~\\
$\mathcal{S} = \left\lbrace \ansbox[4.5cm]{(\texttt{left},\,\texttt{low})}, \ansbox[4.5cm]{(\texttt{left},\,\texttt{high})}, \ansbox[4.5cm]{(\texttt{centre},\,\texttt{low})}, \right.$\\[2mm]
\hfill$\left.\ansbox[4.5cm]{(\texttt{centre},\,\texttt{high})},\ansbox[4.5cm]{(\texttt{right},\,\texttt{low})},\ansbox[4.5cm]{(\texttt{right},\,\texttt{high})}\right\rbrace$
\vspace{80pt}

\item \points[3] Write out the full transition function $p(s' \mid s, a)$ for the state
$s = (\texttt{centre}, \texttt{low})$ under each of the three actions.
Present your answer as a table.

\renewcommand{\arraystretch}{1.9}
\begin{table}[H]
\centering
\begin{tabular}{|l|l|c|}
\hline
Action & Next state $s'$ & $p(s' \mid s, a)$ \\\hline
\multirow{2}{*}{\texttt{move}}
  & $(\texttt{right},\,\texttt{low})$  & \ansnt{$0.5$} \\\cline{2-3}
  & $(\texttt{centre},\,\texttt{low})$ & \ansnt{$0.5$} \\\hline
\texttt{stay}
  & $(\texttt{centre},\,\texttt{low})$ & \ansnt{$1$} \\\hline
\texttt{charge}
  & \multicolumn{2}{c|}{\textit{not available}} \\\hline
\end{tabular}
\end{table}

\item \points[2] A student proposes the following policy:
\begin{quote}
  \emph{``Always \texttt{move}, regardless of battery level or position.''}
\end{quote}
Give a concrete argument (using the reward structure and transition probabilities) for why
this policy may \emph{not} be optimal.
\begin{solution}
For marking: just look for that the student stated that battery levels may change things.
\\~\\
The policy ``always \texttt{move}'' may not be optimal because it ignores the battery level, which directly affects transition probabilities. For instance, from $(\texttt{left},\texttt{low})$, a \texttt{move} succeeds with probability $0.5$, whereas from $(\texttt{left},\texttt{high})$ it succeeds with probability $0.9$. Since the goal is to maximize expected total reward, a policy that takes battery level into account can outperform one that always chooses \texttt{move}.
\end{solution}
% \vspace{80pt}

\item \points[2] A second student claims:
\begin{quote}
  \emph{``The \texttt{charge} action should never be taken, because it always costs $-2$
  and never moves the rover closer to the goal.''}
\end{quote}
Give a concrete counter-argument showing when \texttt{Charge} \emph{can} be
part of an optimal policy.  Your answer should reference specific states and expected
rewards.
% \vspace{80pt}
\begin{solution}
For marking: look that they tried to compute the q values of change vs. moving correctly.
\\~\\
For the last part, under the model exactly as written, the student’s claim is actually consistent with the numbers: from $(\texttt{left},\texttt{low})$, charging gives immediate reward $-2$ and then moves to $(\texttt{left},\texttt{high})$, while moving immediately already has positive success probability and avoids the extra charging cost. In fact,
\[
V(\texttt{left},\texttt{low};\texttt{move-only}) = 7,
\qquad
Q((\texttt{left},\texttt{low}),\texttt{charge}) = -2 + V(\texttt{left},\texttt{high}) \approx 6.78,
\]
so \texttt{charge} is worse here. Thus, with the given rewards and transition probabilities, \texttt{charge} does not appear to be part of an optimal policy.
\\~\\
Details:
Assume that once the rover first reaches \texttt{right}, the episode terminates.

First compute the value of always choosing \texttt{move} from
$(\texttt{centre},\texttt{low})$:
\[
V(\texttt{centre},\texttt{low})
= 0.5(10) + 0.5\bigl(-1 + V(\texttt{centre},\texttt{low})\bigr).
\]
So
\[
V(\texttt{centre},\texttt{low}) = 5 + 0.5V(\texttt{centre},\texttt{low}) - 0.5,
\]
\[
0.5V(\texttt{centre},\texttt{low}) = 4.5,
\]
\[
V(\texttt{centre},\texttt{low}) = 9.
\]

Now compute the value of always choosing \texttt{move} from
$(\texttt{left},\texttt{low})$:
\[
V(\texttt{left},\texttt{low})
= 0.5\bigl(-1 + V(\texttt{centre},\texttt{low})\bigr)
+ 0.5\bigl(-1 + V(\texttt{left},\texttt{low})\bigr).
\]
Substitute $V(\texttt{centre},\texttt{low})=9$:
\[
V(\texttt{left},\texttt{low})
= 0.5(8) + 0.5\bigl(-1 + V(\texttt{left},\texttt{low})\bigr)
= 4 - 0.5 + 0.5V(\texttt{left},\texttt{low}).
\]
Hence
\[
0.5V(\texttt{left},\texttt{low}) = 3.5,
\qquad
V(\texttt{left},\texttt{low}) = 7.
\]

Next compute the value of always choosing \texttt{move} from
$(\texttt{centre},\texttt{high})$:
\[
V(\texttt{centre},\texttt{high})
= 0.9(10) + 0.1\bigl(-1 + V(\texttt{centre},\texttt{high})\bigr).
\]
So
\[
V(\texttt{centre},\texttt{high})
= 9 - 0.1 + 0.1V(\texttt{centre},\texttt{high}),
\]
\[
0.9V(\texttt{centre},\texttt{high}) = 8.9,
\]
\[
V(\texttt{centre},\texttt{high}) = \frac{8.9}{0.9} = \frac{89}{9}\approx 9.89.
\]

Then from $(\texttt{left},\texttt{high})$:
\[
V(\texttt{left},\texttt{high})
= 0.9\bigl(-1 + V(\texttt{centre},\texttt{high})\bigr)
+ 0.1\bigl(-1 + V(\texttt{left},\texttt{high})\bigr).
\]
Substitute $V(\texttt{centre},\texttt{high})=\frac{89}{9}$:
\[
V(\texttt{left},\texttt{high})
= 0.9\left(-1+\frac{89}{9}\right)
+0.1\bigl(-1+V(\texttt{left},\texttt{high})\bigr)
= 8.0-0.1+0.1V(\texttt{left},\texttt{high}).
\]
Hence
\[
0.9V(\texttt{left},\texttt{high}) = 7.9,
\]
\[
V(\texttt{left},\texttt{high}) = \frac{7.9}{0.9} = \frac{79}{9}\approx 8.78.
\]

Therefore, if we charge first at $(\texttt{left},\texttt{low})$ and then always move,
\[
Q\bigl((\texttt{left},\texttt{low}),\texttt{charge}\bigr)
= -2 + V(\texttt{left},\texttt{high})
= -2 + \frac{79}{9}
= \frac{61}{9}\approx 6.78.
\]

So:
\[
V(\texttt{left},\texttt{low};\texttt{move-only}) = 7,
\qquad
Q((\texttt{left},\texttt{low}),\texttt{charge}) = \frac{61}{9}\approx 6.78.
\]
Thus charging is worse under the model
\end{solution}

\end{enumerate}